\section{Đối tượng cần giải thích}
\label{sec:ch02-doi-tuong-can-giai-thich}

\index{mô hình ngôn ngữ lớn}%
Chương trước đã nêu rằng \tn{độ trung thực}{faithfulness} luôn phải gắn với một hành vi và một
loại can thiệp xác định. Với \tn{mô hình ngôn ngữ lớn}{large language model},
hành vi dễ thấy nhất là sinh token kế tiếp: sau một dãy token đầu vào, mô
hình trả một phân phối xác suất trên token kế tiếp, chọn một token từ phân
phối ấy, rồi lặp lại quá trình này để sinh câu trả lời.
Tuy nhiên, người dùng thường không đặt một token rời rạc vào mô hình. Họ đưa
vào một \tn{lời nhắc}{prompt} có chỉ dẫn, dữ kiện, ví dụ, quy ước hội thoại
và đôi khi cả những ràng buộc về định dạng đầu ra.

Vì vậy, một câu trả lời có ít nhất ba lớp cần phân biệt. Lớp thứ nhất là đầu
vào thực tế, gồm các token và thứ tự của chúng. Lớp thứ hai là trạng thái tính
toán do kiến trúc transformer tạo ra khi các token tương tác với nhau. Lớp thứ
ba là hành vi quan sát được, chẳng hạn câu trả lời có làm theo chỉ dẫn hay có
trình bày các bước trung gian hay không. Một lời giải thích ở một lớp không tự
chuyển sang lời giải thích ở lớp khác.

Sự phân biệt này đặc biệt cần thiết khi đọc một giao diện hội thoại. Giao diện
có thể cho thấy một câu trả lời mạch lạc và một chuỗi bước nghe hợp lý, nhưng
nó không hiển thị trực tiếp trọng số, biểu diễn trung gian hay quan hệ nhân quả
giữa chúng với token đầu ra. Do đó, khi các phương pháp ở Phần~II che một phần đầu vào
hoặc thay đổi một \tn{đặc trưng}{feature}, chúng đang kiểm tra hành vi đầu ra dưới một can
thiệp đã chọn, chứ chưa mở được toàn bộ cơ chế bên trong.

Một điều cần nói trước, vì nó quyết định cách đọc cả Phần~II. Các phương pháp
mà Phần~II trình bày đều cần một đầu ra là \emph{một số}, và cần một đơn vị để
can thiệp lên. Với dữ liệu bảng thì cả hai đều sẵn có; với ảnh và văn bản thì
chúng phải được chọn, và lựa chọn ấy là một phần của lời giải thích chứ không
phải chi tiết kỹ thuật. Bảng~\ref{tab:ch02-ba-loai-nhiem-vu} viết ra cả ba
trường hợp.

\begin{table}[htbp]
  \centering
  \footnotesize
  \caption{Ba loại nhiệm vụ mà cuốn sách đụng tới, và cái mà một phương pháp
  attribution cần có ở mỗi loại. Cột giữa là đầu ra vô hướng được giải thích:
  không phương pháp nào ở Phần~II giải thích trực tiếp một câu trả lời bằng văn
  bản.}
  \label{tab:ch02-ba-loai-nhiem-vu}
  \begin{tabular}{@{}p{2cm}p{3cm}p{4cm}p{3cm}@{}}
    \toprule
    Loại dữ liệu & Đầu vào & Đầu ra vô hướng được giải thích & Đơn vị bị can
    thiệp \\
    \midrule
    Bảng & Một vector các đặc trưng & Giá trị dự đoán, hoặc xác suất của một lớp
    & Một cột \\
    \addlinespace
    Ảnh & Một lưới điểm ảnh & Điểm số của lớp được dự đoán & Một điểm ảnh, hoặc
    một mảng điểm ảnh \\
    \addlinespace
    Văn bản & Một dãy token & Xác suất của token hoặc của nhãn đang xét & Một
    token, hoặc một đoạn được đánh dấu \\
    \bottomrule
  \end{tabular}
\end{table}

Với một mô hình ngôn ngữ, cột thứ ba là chỗ phải dừng lại. Câu trả lời sinh ra
là một dãy token, không phải một số, nên muốn dùng LIME hay SHAP trên nó thì
phải chọn lấy một đại lượng vô hướng: xác suất mà mô hình gán cho một token cụ
thể, hoặc cho một nhãn trong một nhiệm vụ phân loại đã đặt ra. Chọn đại lượng
nào là chọn hành vi nào đang được giải thích, tức đúng vế thứ nhất mà
định nghĩa~\ref{def:ch01-do-trung-thuc} đòi phải nêu rõ.
